% !tp-example Basic usage
\documentclass{article}
\begin{document}
\wcTpAnimation[in=shakeX,trigger=click,autoplay,duration=1200ms,iterations=3]{
  \wcTpBox{This box shakes three times on load. Click to replay.}
}
\end{document}

% !tp-example Javascript API
\documentclass{article}
\begin{document}
\wcScript[type=module]|
import "/components/animation/animation.js";
import "/components/box/box.js";

const example = document.createDocumentFragment();

const element1 = document.createElement("tp-animation");
element1.setAttribute("in", "shakeX");
element1.setAttribute("trigger", "click");
element1.setAttribute("autoplay", "");
element1.setAttribute("duration", "1200ms");
element1.setAttribute("iterations", "3");
element1.append(document.createTextNode("\n  "));

const element2 = document.createElement("tp-box");
element2.append(document.createTextNode("This box shakes three times on load. Click to replay."));
element1.append(element2);
element1.append(document.createTextNode("\n"));
example.append(element1);

document.body.append(example);
|
\end{document}

% !tp-example Attention seekers
\documentclass{article}
\begin{document}
Click each box to compare three ways of drawing attention.


\wcTpStack{
  \wcTpAnimation[in=bounce,trigger=click,duration=1200ms]{
    \wcTpBox{Bounce: a vertical rebound.}
  }
  \wcTpAnimation[in=shakeX,trigger=click,duration=1200ms]{
    \wcTpBox{Shake: a horizontal oscillation.}
  }
  \wcTpAnimation[in=pulse,trigger=click,duration=1200ms,iterations=3]{
    \wcTpBox{Pulse: a gentle change of scale.}
  }
}
\end{document}

% !tp-example Directional entrances
\documentclass{article}
\begin{document}
Click each box to compare a sliding entrance with a bouncing entrance.


\wcTpStack{
  \wcTpAnimation[in=slideInLeft,trigger=click,duration=1200ms]{
    \wcTpBox{Slide in from the left.}
  }
  \wcTpAnimation[in=bounceInDown,trigger=click,duration=1200ms]{
    \wcTpBox{Bounce in from above.}
  }
}
\end{document}

% !tp-example Zoom and rotation
\documentclass{article}
\begin{document}
Click a box to play its entrance again.


\wcTpStack{
  \wcTpAnimation[in=zoomIn,trigger=click,duration=1200ms]{
    \wcTpBox{Zoom in: grow from a smaller size.}
  }
  \wcTpAnimation[in=rotateIn,trigger=click,duration=1200ms]{
    \wcTpBox{Rotate in: combine rotation and fading.}
  }
}
\end{document}

% !tp-example Three-dimensional flip
\documentclass{article}
\begin{document}
Hover over the box, or focus it with the keyboard, to flip it around its vertical axis. Leave and return to replay.


\wcTpAnimation[in=flipInY,trigger=hover,duration=1500ms]{
  \wcTpBox{Flip in around the vertical axis.}
}
\end{document}

% !tp-example Entrance and exit
\documentclass{article}
\begin{document}
Use Enter or Exit, then Pause during the movement. Resume continues from the same position; Reset restores the box.


\wcTpCluster{
  \wcTpButton[id=animation-enter]{Enter}
  \wcTpButton[id=animation-exit]{Exit}
  \wcTpButton[id=animation-pause]{Pause}
  \wcTpButton[id=animation-resume]{Resume}
  \wcTpButton[id=animation-reset]{Reset}
}

\wcTpAnimation[id=controlled-animation,in=fadeInLeft,out=fadeOutRight,trigger=manual,duration=3000ms]{
  \wcTpBox{Enter from the left; exit to the right.}
}

\wcScript[type=module]|
await customElements.whenDefined("tp-animation");
const animation = document.getElementById("controlled-animation");
document.getElementById("animation-enter").addEventListener("click", () => animation.playIn());
document.getElementById("animation-exit").addEventListener("click", () => animation.playOut());
document.getElementById("animation-pause").addEventListener("click", () => animation.pause());
document.getElementById("animation-resume").addEventListener("click", () => animation.resume());
document.getElementById("animation-reset").addEventListener("click", () => animation.cancel());
|
\end{document}
